\begin{abstract}
Current evaluations of agentic AI systems primarily rely on benchmarks. However, the static nature of benchmarks makes them uninformative about the value of auditing interventions at deployment time.
% These benchmarks consist of many long-horizon tasks, along with the environments, tools, and context available to agents, as well as criteria for evaluating their current trajectories during task execution.
% While a valuable source of data for characterizing agent behavior, benchmarks have significant limitations in guiding the use of agentic AI in real-world settings.
% Current evaluations of agentic AI typically rely on benchmarks that includes long-horizon tasks, interactive environments, available tools, and criteria for assessing agent trajectories. While these benchmarks provide valuable insights into agent behavior, they offer limited guidance on how interventions during execution might affect downstream behavior and final task outcomes.
% In real world use of agentic AI, humans are often determining within a trajectory whether or not to intervene, e.g., by verifying AI agent output, or by steering the agent to a different behavior.  Whether or not to apply such an intervention is a fundamentally causal inferential question: the user is ascertaining whether the causal effect of such an intervention is beneficial on the final outcome.  This causal effect is particularly challenging because it is a dynamic target: the impact of an intervention depends on the current state of context and environment, and the future behavior of the agentic policy.  Benchmarks as currently constituted are incapable of providing any insight into these causal effects.
For reliable deployment, we often need to determine {\em within an agent's trajectory} whether or not to intervene, e.g., by verifying the agent's current output or steering it toward different behavior, using either human review or more costly models or time-consuming agentic supervision.
The value of such interventions is fundamentally a {\em causal inference} question, asking whether the {\em causal effect} of such an intervention is beneficial for the final outcome.
Estimating this causal effect is particularly challenging due to the dynamic nature of the problem: the impact of an intervention depends on the current context and environment, as well as the agent's future behavior.
% Benchmarks, as currently designed, are incapable of providing insight into these causal effects.
We develop the first methodologically grounded framework for evaluating the causal effects of randomized interventions within agent trajectories.
By leveraging a relatively small dataset of agent trajectories collected under randomized interventions such as verification or screening, we can estimate the dynamic causal effect of a rich class of (contextual) intervention policies. The structural parameters that we estimate can also be used to construct safety leaderboards that capture the risk of various agentic actions, which can be used by practitioners for audit prioritization.
% Remarkably, despite the simplicity of the randomization scheme, we show that these data can be used to credibly estimate the causal effects of a rich class of intervention policies, including those that depend on the agent's observable history.
Applying our methods to $\tau$-bench Retail, ToolSandbox, and ImpossibleBench \citep{yao2025tau,barres2025tau2,lu2025toolsandbox,zhong2025impossiblebench}, we find that intervention effects vary substantially across agent actions and states.
More broadly, our work suggests a rich research agenda at the intersection of causal inference and agentic evaluation, which we believe is essential for the reliable and safe deployment of AI agents in the real world.\footnote{We used ChatGPT, Codex, Claude, and Claude Code to assist with writing, implementing experiments, and analyzing results. The authors take responsibility for the manuscript and results.}


% In this paper, we undertake the first methodologically grounded causal evaluation of such interventions for agentic AI.  Our key insight is that by leveraging a relatively small dataset consisting of agentic trajectories together with randomized application of interventions such as verification or screening, we can estimate the dynamic causal effect of the intervention.  Remarkably, despite simple randomization applied to generate this data, we show that such data can be used to credibly estimate the treatment effect of a very rich class of intervention policies (essentially, any policy that depends only on observable history).  Application of our methods to examples from A, B, C shows D.  More broadly, our paper suggests a rich research agenda at the intersection of causal inference and agentic evaluation, that is essential if we hope to reliably and safely deploy agents in the real world.

% TODO(AppWorld removal): Commented draft mentions/cites AppWorld; remove or update if reused.
% AI agents increasingly interact with external environments through sequences of decisions and tool calls, yet existing benchmarks primarily evaluate terminal task success or safety violations. Such evaluations provide limited insight into how interventions at intermediate decisions affect downstream agent behavior and final outcomes. In this paper, we adapt micro-randomized trials and methods from sequential causal inference to evaluate interventions within agent trajectories. By randomizing intervention assignments at eligible decision points, a single collection of trajectories can support causal evaluation across different agent actions and states. We formalize three causal estimands that distinguish the effects of individual interventions under different continuation rules from the gains of selective intervention policies. We apply augmented inverse probability weighting (AIPW) and structural nested mean models (SNMM) to estimate these effects with uncertainty quantification. We evaluate this approach in three agent environments: tool-call review in $\tau$-bench Retail, code-block advice in AppWorld, and adversarial steering in ImpossibleBench. Our results reveal heterogeneous intervention effects, with targeted assistance improving task success and adversarial feedback increasing safety violations in certain agent states. 
% These findings illustrate the potential of randomized interventions for causal evaluation of agent capabilities and safety, while highlighting the statistical challenges of estimating intervention-policy effects over long trajectories.
\end{abstract}


